\section*{Chapitre \ref{ch:fm:managing-a-book-through-drawdown} --- Gérer un livre pendant un drawdown}

\begin{solution}{exo:fm:managing-a-book-through-drawdown:1}
0.8 de la volatilité annuelle à 10~\%~; 0.27 à 30~\%.
\end{solution}

\begin{solution}{exo:fm:managing-a-book-through-drawdown:2}
$A=2\,520$ jours~: $2\ln2\,520/(1-0)^2=15.7$ ans~; jusqu'à $-1$~: $15.7/4=3.9$ ans.
\end{solution}

\begin{solution}{exo:fm:managing-a-book-through-drawdown:3}
$9.79\times10=98~\%$ des stratégies saines sur dix ans~; 66~\% des stratégies cassées avant leur rupture.
\end{solution}

\begin{solution}{exo:fm:managing-a-book-through-drawdown:4}
Plus élevée~: la probabilité a priori devient $0.30+0.70/1\,260=0.3006$, et le rendement favorise l'état cassé par un rapport de vraisemblance de
seulement $\exp(\tfrac12(S/\sqrt{252})^2)=1.002$, de sorte que la probabilité a posteriori vaut 0.3010. Le rendement d'un jour porte un signal de $S/\sqrt{252}=0.063$ écart type.
\end{solution}

\begin{solution}{exo:fm:managing-a-book-through-drawdown:5}
Son P\&L espéré est nul après la rupture, donc l'arrêter ne fait perdre ni ne rapporte rien en espérance. Cela libère le capital et le risque qu'elle
consomme~: la \omterm{def:fm:limits-and-capital-allocation:raroc}{charge en capital} pendant son fonctionnement, et la variance qu'elle ajoute à la firme.
\end{solution}

\begin{solution}{exo:fm:managing-a-book-through-drawdown:6}
Non~: 12~\% est inférieur au drawdown maximal médian sur dix ans d'une stratégie saine (16.8~\%)~; des drawdowns de cette profondeur sont normaux et
en disent peu sur une rupture.
\end{solution}

\begin{solution}{exo:fm:managing-a-book-through-drawdown:7}
La règle a posteriori détecte les stratégies cassées après 1.3 an en médiane (fausses alertes~: 2.24 pour cent ans)~; les valeurs pour une population
moitié-moitié sont 17.5 (stop-loss), 23.2 (échelle), 31.3 (\omterm{def:fm:managing-a-book-through-drawdown:rules}{stop temporel}) et 39.3 (a posteriori), contre 10.1 sans règle. Lorsqu'une rupture se traduit par
des pertes, chaque règle aide, et la règle a posteriori le plus.
\end{solution}

\begin{solution}{exo:fm:managing-a-book-through-drawdown:8}
Le backtest compte les pertes évitées et non les reprises manquées~: le stop se déclenche sur des stratégies saines lors de drawdowns ordinaires (98~\% d'entre
elles sur dix ans dans le modèle du chapitre). Compter les fausses alertes et le P\&L perdu après elles, sur des stratégies dont on sait qu'elles sont saines, avant de l'adopter.
\end{solution}

\begin{solution}{pb:fm:managing-a-book-through-drawdown:1}
\begin{enumerate}
\item Voir la \cref{def:fm:managing-a-book-through-drawdown:rules}.
\item 16.8~\% et 25.4~\% du capital.
\item 0.8 et 0.4 de la volatilité annuelle.
\item Les drawdowns d'une stratégie saine varient avec sa volatilité~; un seuil exprimé en pourcentage n'a pas la même signification pour des stratégies différentes.
\item Voir la \cref{prop:fm:managing-a-book-through-drawdown:posterior}.
\item 0.24 pour les stratégies saines et 0.39 pour les stratégies cassées (seulement 13~\% des trajectoires cassées atteignent d'abord 8~\% après leur rupture).
\item Délai $2\ln A/(S-S_0)^2$ années~: 15.7 et 3.9.
\item Taux de capture et d'exécution, mesures d'encombrement, changements de structure de marché~: des observations qui évoluent avec l'edge lui-même, et non avec le
bruit de la stratégie.
\item Fausses alertes~: 9.79, 2.09, 6.19 et 2.84 pour cent années-stratégie~; ruptures détectées~: 34~\%, 60~\%, 80~\% et 95~\%.
\item 0.72, 2.79, 2.00 et 2.68 ans.
\item Vivantes 21.7, 39.5, 54.6, 71.4~; cassées 15.0, 12.8, 19.5, 18.2~; moitié-moitié 18.3, 26.2, 37.1, 44.8~; sans règle 79.0, 11.1, 45.1.
\item Elle arrête presque toutes les stratégies saines lors d'un drawdown ordinaire, et la plupart des stratégies cassées avant leur rupture.
\item Arrêter une stratégie cassée au Sharpe 0 ne fait économiser que sa \omterm{def:fm:limits-and-capital-allocation:raroc}{charge en capital}~; elle en économise davantage lorsque les stratégies cassées perdent de l'argent (exercice 7).
\item Voir la \cref{def:fm:managing-a-book-through-drawdown:restart}~: il faut faire du paper trading de la stratégie arrêtée et la redémarrer à taille réduite lorsque son P\&L
en paper trading se redresse.
\item Réexaminer la thèse et les preuves de la stratégie comme si elle était nouvelle~; cela utilise une information autre que le P\&L.
\item Elles poussent les gérants à jouer leur va-tout ou à se figer~; les échelles et les \omterm{def:fm:managing-a-book-through-drawdown:restart}{règles de redémarrage} atténuent ces deux effets.
\item Comparer les deux drawdowns en unités de volatilité et avec les distributions saines~; examiner des preuves autres que le P\&L~; ne réduire la taille que si elles
indiquent une rupture.
\item La règle a posteriori avec un taux de hasard tiré de l'expérience de la firme, une échelle pour la taille, et une \omterm{def:fm:managing-a-book-through-drawdown:restart}{règle de redémarrage}~; aucun stop fondé sur la seule profondeur.
\item Probabilité a posteriori de 0.24 (saine) et 0.39 (cassée) au drawdown de 8~\%~; 15.7 ans pour détecter une chute de 1 à 0 avec une fausse alarme tous les dix ans.
\item On vous jugera sur les preuves que votre edge est intact, dont votre drawdown n'est qu'un élément faible~; votre taille baissera progressivement à mesure que
les preuves contraires s'accumuleront, et un arrêt ne sera pas une fin si les preuves se redressent.
\end{enumerate}
\end{solution}

\begin{solution}{iq:fm:managing-a-book-through-drawdown:1}
Pas en soi~: 15~\% représentent 1.5 année de volatilité, dans la plage normale du drawdown maximal sur dix ans d'une stratégie de Sharpe 1 (médiane d'environ 17~\%).
\iqlookfor{unités de volatilité et distribution saine.}
\end{solution}

\begin{solution}{iq:fm:managing-a-book-through-drawdown:2}
Un stop met fin au trading à un seuil~; une échelle réduit la taille par paliers. L'échelle produit moins de fausses alertes et réagit avant le stop.
\iqlookfor{réponse graduée.}
\end{solution}

\begin{solution}{iq:fm:managing-a-book-through-drawdown:3}
KL par jour $=(S-S_0)^2/(2\times252)=1/504$~; délai $\approx\ln A\times504$ jours, soit environ 16 ans pour $A=2\,520$.
\iqlookfor{information par jour et logarithme de la période entre fausses alarmes.}
\end{solution}

\begin{solution}{iq:fm:managing-a-book-through-drawdown:4}
Prédire $\tilde p=p+(1-p)h$, puis appliquer Bayes avec les deux vraisemblances gaussiennes (\cref{prop:fm:managing-a-book-through-drawdown:posterior}).
\iqlookfor{les étapes de prédiction et de mise à jour.}
\end{solution}

\begin{solution}{iq:fm:managing-a-book-through-drawdown:5}
À quelle fréquence elle aurait arrêté des stratégies qui se sont ensuite redressées, et ce que valaient ces redressements.
\iqlookfor{fausses alertes et reprises manquées.}
\end{solution}

\begin{solution}{iq:fm:managing-a-book-through-drawdown:6}
Taux de capture, d'exécution et de réussite, décroissance du signal, encombrement et flux~: ils mesurent l'edge directement, avec beaucoup moins de bruit par jour que le P\&L.
\iqlookfor{observations à meilleur rapport signal sur bruit.}
\end{solution}
